% paper/sections/08_1_bf_seven_principles.tex
%
% §7.3  Balanced-Force 七原則
%       P-1 位置一致 / P-2 PPE-corrector 共有 / P-3 geometry 経路統一 /
%       P-4 β_f 統一 / P-5 曲率への CCD 非適用 /
%       P-6 界面ジャンプ補正 / P-7 BF sub-system 分離
%
% ═══════════════════════════════════════════════════════════════════════════════
\subsection{Balanced-Force 七原則}
\label{sec:bf_seven_principles}
% ═══════════════════════════════════════════════════════════════════════════════

§\ref{sec:bf_failure_modes} の 5 失敗モードを回避するため，
本節では 7 原則 P-1..P-7 を与える．
\textbf{BF 整合性の核は P-1（位置一致）/ P-2（補正 $G_h$ 共有）/ P-4（$\betaf$ 統一）/
P-6（界面ジャンプ補正）の 4 原則}が担い，これらが欠けると寄生流れが
$\Ord{h^0}$ で残留する．残る P-3（geometry 経路統一）/ P-5（曲率非 CCD 化）/
P-7（sub-system 分離）は\textbf{次数改善および予防的役割}を持つ．
\textbf{特に P-5 の status は適用形態に依存}：対称フィルタ適用時は
P-1 維持の必要条件として連動する（§\ref{sec:bf_p5_curvature}），
非対称・無フィルタ時は次数改善の予防策に留まる．主解消は依然 4 原則 (P-1/P-2/P-4/P-6) で完結する
（詳細は §\ref{sec:bf_p5_curvature} および §\ref{sec:bf_failure_principle_map} の F$\times$P 表参照）．

% -------------------------------------------------------------------------------
\subsubsection{\texorpdfstring{P-1：$G_h p$ と $f_{\sigma,h}$ の同一位置評価}{P-1：Gh p と f sigma h の同一位置評価}}
\label{sec:bf_p1_location}
% -------------------------------------------------------------------------------

\textbf{コロケート配置の面演算子 FCCD（\S\ref{sec:fccd_def}）で実装する場合}，
速度補正は面値 $u_f^{n+1}$ を経由して節点に戻す 2 段設計となる：
\begin{equation}
  u_f^{n+1} = u_f^* + \Delta t_{\mathrm{proj}}\Bigl(
    -\betaf(G_h^\text{bf} p)_f + (f_\sigma)_f + \cdots
  \Bigr),
  \label{eq:bf_p1_corrector}
\end{equation}
$(G_h^\text{bf} p)_f$ と $(f_\sigma)_f$ は\textbf{同じ面で評価}されねばならない．
節点 CCD 勾配を面に補間し，$f_\sigma$ を別の補間で構築する
\textbf{「位置を合わせれば十分」は不正解}．
補間演算子自身が互いに随伴整合（adjoint-compatible）でなければ
$(G_h p)_f - (f_\sigma)_f$ は相殺しない．
P-1 から P-7 まで本節全体を通じ，「面」は FCCD で定義される共通面ジェット
位置（\S\ref{sec:face_jet_def}）を指す．

% -------------------------------------------------------------------------------
\subsubsection{\texorpdfstring{P-2：随伴 PPE による F-2・F-3 の同時解消：$D_h^\text{bf} = -(G_h^\text{bf})^\ast$}{P-2：随伴 PPE による F-2・F-3 の同時解消}}
\label{sec:bf_p2_adjoint}
% -------------------------------------------------------------------------------

PPE を\textbf{随伴ペア}から組み立てる：
\begin{equation}
  D_h^\text{bf}\Bigl(\betaf\,G_h^\text{bf}\,p\Bigr)
    = \frac{1}{\Delta t_{\mathrm{proj}}}\,D_h^\text{bf}\,\mathbf{u}^*,
  \qquad D_h^\text{bf} = -(G_h^\text{bf})^\ast.
  \label{eq:bf_p2_ppe}
\end{equation}
これにより離散内積 $\langle\cdot,\cdot\rangle_h$ 下で
\begin{equation}
  \langle p,\,A_h p\rangle_h
    = \langle G_h^\text{bf} p,\,\betaf\,G_h^\text{bf} p\rangle_h
    \ge 0,
  \label{eq:bf_p2_spd}
\end{equation}
零空間（定数モード）除去後に $A_h$ は SPD で CG が主解法，
\textbf{補正側 $G_h$ = PPE 側 $G_h^\text{bf}$} が保証される．

\noindent\textbf{rising bubble への含意}：
PPE RHS の発散だけを修正しても，補正側が別の $G_h$ を参照する限り
F-2 は解消しない．$G_h^\text{bf}$ を両経路で共有することが必要．

\noindent\textbf{F-2 と F-3 の独立性}：
F-3（$D_h\neq -G_h^\ast$；随伴ペア違反）と
F-2（PPE 内 $G_h\neq$ 補正 $G_h^\text{bf}$；経路違反）は\textbf{独立の失敗モード}だが，
P-2 の随伴構築 $D_h^\text{bf}=-(G_h^\text{bf})^\ast$ は両者を同時に解消する：
$G_h^\text{bf}$ を固定して $D_h^\text{bf}$ をその随伴として導出するため，
PPE と補正の両経路で同一演算子が現れ，かつ随伴整合も自動的に成立する．

% -------------------------------------------------------------------------------
\subsubsection{\texorpdfstring{P-3：$f_\sigma$ を $[p]_\Gamma$ と同じ離散幾何に通す}{P-3：f sigma を圧力ジャンプと同じ離散幾何に通す}}
\label{sec:bf_p3_geometry}
% -------------------------------------------------------------------------------

CSF 形式の $f_\sigma$ を，仮想毛管圧 $p_\sigma$ から
$G_h^\text{bf} p_\sigma$ として構成できるよう設計すれば，
静止平衡は
\begin{equation}
  -G_h^\text{bf} p + G_h^\text{bf} p_\sigma
    = G_h^\text{bf}(p_\sigma - p) \approx \mathbf{0}
  \label{eq:bf_p3_cancel}
\end{equation}
が\textbf{PPE 解の精度}により成立する．
独立に離散化された 2 ベクトルのキャンセルを期待する旧来設計と本質的に異なる．

% -------------------------------------------------------------------------------
\subsubsection{\texorpdfstring{P-4：$\betaf$ を 3 経路で統一}{P-4：beta f を 3 経路で統一}}
\label{sec:bf_p4_beta}
% -------------------------------------------------------------------------------

変密度 PPE $L_h p = -D_h^\text{bf}(\betaf\,G_h^\text{bf} p)$ において
$\betaf = (1/\rho)_f$ は以下の 3 経路\textbf{すべてで同一}：
\begin{enumerate}[leftmargin=2em,noitemsep]
  \item PPE 演算子 $L_h$ 自身，
  \item 速度補正 $u_f^{n+1} = u_f^* - \Delta t_{\mathrm{proj}}\,\betaf\,(G_h^\text{bf} p)_f$，
  \item 界面張力正規化 $f_\sigma \propto \betaf\,\sigma\kappa\,(\nabla\psi)_f$（CSF
        標準式 $\sigma\kappa\nabla\psi$ に $\betaf$ 因子を付加した\textbf{変密度
        corrector 離散化の特殊形}；等密度では $\betaf$ が定数で標準式に還元する）．
\end{enumerate}
PPE に調和平均，補正に算術平均という「混用」は
$\rho_l/\rho_g \gg 1$ で BF 残差を界面で $\Ord{1}$ とする．
\textbf{本稿の選択}：$\betaf^\text{harm}=2/(\rho_L+\rho_R)$ を 3 経路共通で採用．

\smallskip\noindent\textbf{命名と等価性（混乱回避）}：
$a\equiv 1/\rho$ と置くと，FVM フェイス係数は
\textbf{$a$ の調和平均}
\begin{equation}
  \betaf^\text{harm}
  = a_f^\text{harm}
  = \frac{2a_La_R}{a_L+a_R}
  = \frac{2}{\rho_L+\rho_R}
  \label{eq:bf_beta_harmonic}
\end{equation}
である（§\ref{app:fvm_face_coeff}）．
これは 2 値界面では\textbf{密度の算術平均の逆数}
$[(\rho_L+\rho_R)/2]^{-1}$ と同じ値を与える．
一方，\textbf{$\rho$ の調和平均の逆数}
$(\rho_L+\rho_R)/(2\rho_L\rho_R)$ は別物であり，
高密度比では界面係数を誤る典型的な実装ミスである．

% -------------------------------------------------------------------------------
\subsubsection{P-5：曲率に生 CCD を適用しない}
\label{sec:bf_p5_curvature}
% -------------------------------------------------------------------------------

$\kappa = \bnabla\cdot(\bnabla\psi/|\bnabla\psi|)$ は
界面から $\Ord{\varepsilon}$ 幅でしか滑らかでない場を 2 回微分する．
高次 CCD は分母の $|\bnabla\psi|$ 小のノイズを増幅し，
ほぼ定数の bulk を広いステンシルで走査する．

\noindent\textbf{規則}：$\bnabla\psi$ には CCD を用いてよい．
$\kappa$ は\textbf{界面帯に制限した平滑化}または
2 次 FD に落とす．曲率精度は\textbf{演算子位置整合に比べ 2 次的}であり，
BF 残差への影響は 2 次効果に留まる．

% -------------------------------------------------------------------------------
\subsubsection{P-6：界面越えステンシルへのジャンプ補正}
\label{sec:bf_p6_jump}
% -------------------------------------------------------------------------------

CCD は滑らかな場を仮定する．静止液滴で $\Delta p = \sigma\kappa$ が
正しく設定されているのに寄生流れが出る場合，
ほぼ常に\textbf{界面越えステンシルが圧力ジャンプを平滑化}している．
必要な補正：
\begin{itemize}[noitemsep]
  \item \textbf{GFM 面フラックス補正}（Phase 1 実用）：
  \begin{equation}
    (G_h^\text{bf} p)_f^\text{corrected}
      = (G_h^\text{bf} p)_f - \frac{\sigma\kappa}{H_f},
    \label{eq:bf_p6_gfm}
  \end{equation}
  bulk ステンシルは変更せず，ジャンプだけを面フラックスに注入．
  2 次精度；Phase 1 検証と H-01 診断に十分．
  \item \textbf{IIM コンパクト行補正}（Phase 3 研究）：
  FCCD のコンパクト関係 $M_f\mathbf{g} = R_f\mathbf{p}$ に対し
  $[p]_\Gamma,\,[\betaf\partial_n p]_\Gamma$ を Taylor 展開した
  補正ベクトル $c_\Gamma^\text{IIM}$ を加算．
\end{itemize}

% -------------------------------------------------------------------------------
\subsubsection{P-7：BF sub-system と bulk 微分子の分離}
\label{sec:bf_p7_separation}
% -------------------------------------------------------------------------------

\begin{center}
\begin{tabular}{@{}lll@{}}
\hline
タスク & 演算子 & 精度要求 \\
\hline
bulk 移流・拡散（$u,v$）& 節点 CCD / UCCD6 / FCCD & 高次（$\Ord{h^6}$） \\
PPE + 補正 & $G_h^\text{bf},\,D_h^\text{bf}$（面随伴ペア）& 内整合（$\Ord{h^2}$--$\Ord{h^4}$） \\
界面張力 $f_\sigma$ & $G_h^\text{bf}$ と同じ面スロット & 内整合 \\
曲率 $\kappa$ & 界面帯制限平滑化 & 2 次 \\
\hline
\end{tabular}
\end{center}

\noindent\textbf{設計哲学}：
BF sub-system は\textbf{高次である必要はない}．
内部整合性（位置一致・随伴ペア・$\betaf$ 統一）があれば
寄生流れは演算子精度ではなく
\textbf{界面ジャンプ処理の精度}で律速される．
bulk と BF で演算子を分けるのは「妥協」ではなく\textbf{設計原則}．

% -------------------------------------------------------------------------------
\subsubsection{F-1--F-5 × P-1--P-7 対応マップ}
\label{sec:bf_failure_principle_map}
% -------------------------------------------------------------------------------

5 失敗モード（§\ref{sec:bf_failure_modes}）と 7 原則の対応関係：
\begin{center}
\renewcommand{\arraystretch}{1.2}
\begin{tabular}{@{}lcccccccl@{}}
\hline
 & P-1 & P-2 & P-3 & P-4 & P-5 & P-6 & P-7 & 主解消原則 \\
\hline
F-1（位置ズレ）              & $\blacksquare$ & & $\square$ & & & & $\square$ & P-1 \\
F-2（補正 $G_h$ 不一致）      & & $\blacksquare$ & & & & & $\square$ & P-2 \\
F-3（随伴違反）              & & $\blacksquare$ & & & & & & P-2 \\
F-4（$\betaf$ 混用）         & & $\square$ & & $\blacksquare$ & & & & P-4 \\
F-5（界面越えステンシル）\footnote{F-5 で P-6 単独は \(\Ord{h^2}\) を達成する；P-3 と P-5 は次数改善の上乗せ寄与であり主解消ではない．}   & & & $\square$ & & $\square$ & $\blacksquare$ & & P-6 \\
\hline
\end{tabular}

{\footnotesize $\blacksquare$：主解消，$\square$：副次的に寄与}
\end{center}

\noindent P-5（曲率 $\kappa$ の非 CCD 化）と P-3（$f_\sigma$ を $p$ と同幾何）は
F-5（界面越えステンシル平滑化）に対し副次的・間接的に寄与する．
P-7（sub-system 分離）は F-1・F-2 の\textbf{設計レベル}での発生を防ぐ予防策．
\textbf{F-1..F-5 への主解消は P-1, P-2, P-4, P-6 の 4 原則}が担い，
P-3, P-5, P-7 は次数改善・設計予防として\textbf{独立な設計目的}を持つ
（各原則の設計動機は §\ref{sec:bf_p1_location}--\ref{sec:bf_p7_separation} 参照）．
7 原則の動機は単一の「5 モード解消最小集合」ではなく，
主解消 4 原則と副次寄与 3 原則が\textbf{独立な目的群}を網羅する設計である．

% ═══════════════════════════════════════════════════════════════════════════════
